\section{Experimental design}
\label{sec:experiments}

The benchmark treats every method as a feature ranker. The design goal is to
separate ranking quality from downstream model choice while keeping
confirmatory claims tied to a small set of pre-specified endpoints. Broader
summaries are still useful, but they are reported as descriptive robustness
checks rather than as headline evidence.

\subsection{Rank-then-evaluate protocol}
\label{sec:experiments-protocol}

Each method first produces a feature ranking on the training split. We then
train downstream predictors on the top-$k$ features and evaluate them on
held-out data. This rank-then-evaluate protocol isolates the quality of the
ranking itself rather than rewarding methods for a particular internal
training objective. We use 5-fold cross-validation repeated over 5 random
seeds and evaluate $k \in \{5,10,25,50,100,p\}$, clipped to $k \le p$. All
preprocessing is fit on the training portion of each fold only.

\subsection{Primary endpoints and aggregation}
\label{sec:experiments-analysis-plan}

The benchmark produces many numbers across tasks, downstream models, and
values of $k$, so the paper distinguishes \emph{confirmatory} endpoints from
\emph{descriptive} summaries. For classification, the primary real-data
endpoint is balanced accuracy of logistic regression at $k=25$. For
regression, it is $R^2$ of ridge regression at $k=25$. Dataset-level
summaries average across folds and random seeds after fixing the downstream
model and $k$.

Broader summaries that average across multiple downstream models or multiple
values of $k$ are reported separately and labeled as exploratory. They help
show whether a qualitative pattern is stable, but they are not used for exact
headline claims.

For synthetic datasets with known informative sets $S^\star$, the primary
ranking metrics are precision@k, recall@k, and F1@k at
$k \in \{5,10,20,|S^\star|\}$, again clipped to $k \le p$. In null-style
diagnostics we additionally report rejection rates and noise-selection
behavior.

\subsection{Datasets and preprocessing}
\label{sec:experiments-datasets}

We evaluate on curated synthetic suites with known signal structure and on
real tabular datasets spanning moderate-dimensional and $p \gg n$ regimes. The
synthetic suites are designed to isolate distinct failure modes:
high-cardinality null features, nonlinear dependence, correlated blocks,
redundancy, weak signal, and high-dimensional noise injection. The real-data
benchmark is used to assess whether the same qualitative patterns persist
outside controlled designs.

Within each fold, preprocessing is learned on the training portion only and
applied unchanged to the held-out portion. Features are standardized using
training-fold means and variances. Dataset inventories and code-aligned
preprocessing details are collected in the supplement.

\subsection{Compared methods and compute budgets}
\label{sec:experiments-methods}

The benchmark includes three families of rankers.
\emph{Permutation-calibrated filters} turn univariate dependence scores into
permutation p-values. \emph{Embedded methods} derive rankings from trees and
ensemble learners, including conditional-inference forests and common
CART-style baselines. \emph{Wrappers} include recursive elimination and
permutation-importance style procedures. Implementation-aligned definitions of
the selector and splitter statistics used by the conditional-inference models
are given in \cref{app:methods-selectors,app:methods-splitters}.

To reduce optimistic bias from unequal tuning effort, every method is run
under a locked hyperparameter grid and compute budget. Cross-dataset
confirmatory summaries are reported at the method-family level rather than at
the individual configuration-hash level, so small within-method configuration
changes do not dominate the paper story.

\paragraph{Reproducibility.}
Experiment orchestration, analysis, and figure generation scripts live under
\texttt{paper/scripts/}, and generated artifacts are written to
\texttt{paper/results/}. The full benchmark is designed to be reproducible
from the repository.
